% Generated shared reviewed material.
% External preprint input and application proof; no new sieve proof is claimed.
\section{Proof of the physical-row sieve application}\label{phy:details}
\subsection{The precise literature input and the primary generator}
The external input is de Faveri,
\href{https://arxiv.org/abs/2610.04045v1}{\emph{Optimal large sieve for
fixed order characters}, arXiv:2610.04045v1, 2 October 2026},
Theorem~1.1, printed page~2. Input~\eqref{phy:input} is invoked with
the family in its Sections~2.2--2.5; its deep proof is not reproduced
or claimed as a new sieve or Lean theorem here. The primary PDF SHA-256 is the concatenation of these two lines:
\par\noindent\texttt{97720b3cacccbedb9c2014520ad4d128}\par
\noindent\texttt{ae23afd0a781bdeff27dab4e86b7c726}\par
The public citation receipt records the version and text hash. Full
PDF/text reading snapshots are local cache and excluded from public closure.

The ring $\mathcal O=\mathbb Z[\omega]$ is a PID. Its six units map
bijectively onto $(\mathcal O/(3))^\times$, which has six elements.
Consequently every fractional ideal prime to $3$ has a unique generator
$\alpha_{\mathfrak a}\equiv1\pmod3$, interpreted locally at $3$.
This choice is multiplicative, since a product of two such generators
is again $1$ modulo $3$ and generates the product ideal.

Choose the source's fixed modulus $\mathfrak c_0$ sufficiently divisible
at $S$ that its stipulated local sixth-power condition holds, and choose
it divisible by $(3)$. In the source's Section~2.3 choose ordinary ideal
representatives $E_0$ prime to $S$, and set
$m_{E_0}=\alpha_{E_0}$. Global principality makes these allowed
generators after $S$-localization too. The prescribed products then give
$m_E=\alpha_E$ for each representative $E$.
For coprime good ideals $\mathfrak a,\mathfrak m$, the construction gives
\[
 \mathfrak a=(x)E\mathfrak g^6,\qquad
 x\equiv1\pmod{\mathfrak c_0},\qquad
 (\mathfrak g,\mathfrak m)=1,
 \qquad
 \chi_{\mathfrak a}(\mathfrak m)
       =\left(\frac{xm_E}{\mathfrak m}\right)_6.
\]
Choose an ordinary fractional generator $\gamma_{\mathfrak g}$ of
$\mathfrak g$. Ordinary ideal equality implies
\begin{equation}\label{phy:generator-proof}
 \alpha_{\mathfrak a}=\eta x m_E\gamma_{\mathfrak g}^6,
          \qquad\eta\in\mu_6.
\end{equation}
All factors are units at $3$, and
$\gamma_{\mathfrak g}^6\equiv1\pmod3$ because the residue unit
group has order six. The other three displayed generators are $1$
modulo $3$. Hence $\eta\equiv1\pmod3$ and injectivity on $\mu_6$
forces $\eta=1$. Taking the sixth-power residue symbol proves the exact
identification in Proposition~\ref{phy:family}. For a common prime the
source's induced character and the original raw symbol are both zero.
Thus no primitive-conductor transfer can silently restore a nonunit value.

The field contains $\mu_6$; it need not contain $\mu_{12}$ for the
Theorem~1.1 input. Its characters have order dividing six and trivial
infinite type, so physical valuations $2,3,4,5$ are included. They are
not replaced by a squarefree radical or indexed by its norm. The fixed
choice of $S,\mathfrak c_0,E,m_E$ is independent of $D,H,U$ and of all
coefficient masks below. We fix the complex embedding to agree with the
original sextic symbol. This is a deliberate allowed construction, not
an assertion that every noncanonical choice in the source is identical.

\subsection{The complete raw row family and the sixth factor}\label{phy:raw-proof}
At each prime outside $S$, divide its valuation in $(u)$ by six.
The remainder defines the sixth-power-free ideal $\mathfrak a$, and the
quotient defines $\mathfrak r_6$. The quotient may exceed one and may
be positive at a prime of $\mathfrak a$. Thus $\mathfrak r_6$ is neither
assumed squarefree nor made coprime to the first five factors.
With the multiplicative primary generators, the quotient
$u_0=u/(\alpha_{\mathfrak a}\alpha_{\mathfrak r_6}^6)$
is an integral $S$-supported element, including its literal unit.
This decomposition is unique and has the norm relation in~\eqref{phy:raw}.

For a squarefree good column $\mathfrak m$, the sixth-power factor has
symbol one when $(\mathfrak m,\mathfrak r_6)=1$ and zero otherwise.
The exact family identification gives~\eqref{phy:raw}, even when
$\mathfrak a,\mathfrak r_6$ overlap. In particular a prime common to
both factors is still an original raw zero, not two independent masks.
For fixed $u_0,\mathfrak r_6$ set
\[
 c_{\mathfrak m}=A_U(\mathfrak m)\nu(\mathfrak m)
       \chi_{\mathfrak m}(u_0)W(N\mathfrak m/D)
                       \mathbf1_{(\mathfrak m,\mathfrak r_6)=1}.
\]
These coefficients are independent of the remaining row $\mathfrak a$.
They vanish outside the COMPLETE original good squarefree annulus;
putting zero on the extra sixth-power-free indices is permitted.
The actual row restrictions can only decrease this positive norm, and
\begin{equation}\label{phy:coefficient-norm}
 \sum_{\mathfrak m}|c_{\mathfrak m}|^2
       \ll_\epsilon B_W^2D^{1+\epsilon}.
\end{equation}
Indeed $|A_U(\mathfrak m)|\le3^{\omega(\mathfrak m)}$
uniformly in $U$, and there are $O_K(D)$ ideal columns. This is the
original assignment bound; it does not assert
$A_U(dn)=A_U(d)A_U(n)$ or use the rough short-coefficient formula on $K_0$.

The fixed $\nu$ and literal/$S$ phase have modulus one on good columns.
A row-independent unitary Jacobi factor of nonzero infinity type could
also be included in these arbitrary coefficients. This does not turn it
into an order-six row character, supply a theorem for varying infinite
types, or improve the operator norm by a gauge. Arbitrary positive row
weights are not part of this application; they require a true weighted
operator bound in the exact identities of Proposition~\ref{cmr:joint}.

\subsection{The raw global sum and the principal assembly}\label{phy:r6-sums}
For fixed $u_0,\mathfrak r_6$, use~\eqref{phy:input} with
\[
 A=\frac{H}{Nu_0(N\mathfrak r_6)^6},\qquad B=2D.
\]
If $A<1$ there is no remaining ideal row. If $A=1$ the unit row remains;
it is not dropped as a subunit endpoint. Equations~\eqref{phy:raw}
and~\eqref{phy:coefficient-norm}, and the factor $D^{-1}$ in $|I_u|^2$,
give a bound by $B_W^2(HD)^\epsilon\Phi_6(A,D)$.
Now sum all $\mathfrak r_6$ of norm at most $(H/Nu_0)^{1/6}$.
The $A,A^{5/6}D^{1/3},A^{1/3}D^{5/6}$ terms use respectively
\[
 \sum_{\mathfrak r_6}(N\mathfrak r_6)^{-6},\qquad
 \sum_{\mathfrak r_6}(N\mathfrak r_6)^{-5},\qquad
 \sum_{\mathfrak r_6}(N\mathfrak r_6)^{-2},
\]
which converge. The $D$ term instead uses the ideal count
$O_K((H/Nu_0)^{1/6})$. Thus the four resulting terms are
\begin{equation}\label{phy:S-sums}
 \begin{aligned}
 &H(Nu_0)^{-1},\qquad DH^{1/6}(Nu_0)^{-1/6},\\
 &H^{5/6}D^{1/3}(Nu_0)^{-5/6},\qquad
       H^{1/3}D^{5/6}(Nu_0)^{-1/3}.
 \end{aligned}
\end{equation}
Every exponent on the fixed finite $S$-supported semigroup is positive,
so these sums converge geometrically, with the six literal units included.
This proves the raw part of Proposition~\ref{phy:global}, with every
arbitrary sixth-factor overlap still represented by its coefficient mask.

For the residue-subtracted family, use the original formula
\eqref{wc:principal}. Its deleted Euler residue is bounded by a fixed
constant and
\[
 \sum_d\frac{|\beta_U(d)|}{Nd}
   \le\left(\sum_{Na\le U}\frac1{Na}\right)^2
   \ll_K\log^2(2U).
\]
Therefore $|P_u|\ll_{\nu,S}B_W\sqrt D\log^2(2U)$ on a principal row.
Outside $S$, such a row has all valuations divisible by six, since the
original $\nu$ has its defining conductor in $S$. The already established
count is $O_{\nu,S,\epsilon}(H^{1/6}(HD)^\epsilon)$. Consequently
\begin{equation}\label{phy:principal-proof}
 \sum_u|P_u|^2\ll_{\nu,S,\epsilon}
               B_W^2(HD)^\epsilon DH^{1/6}.
\end{equation}
Using $|I-P|^2\le2|I|^2+2|P|^2$ proves the stated $E_u$ bound.
No Artin holomorphy, spectral positivity or zero-location assertion is used.
For $1\le H\le D$, each of the four terms in $\Phi_{\rm all}$ is at most
$DH^{1/6}$, proving~\eqref{phy:dominance}.
The varying-$W_D$ construction of Section~\ref{prf:core} gives a compatible
supremum-norm-class obstruction; it does not make a fixed-$W$ lower bound.

\subsection{The signed $A_6=1$ common-divisor form}\label{phy:divisor-proof}
For a fixed $u_0$ the first-five physical ideal is sixth-power-free.
For the two-pivot block it is exactly $v_1v_2^2$, of norm at most
$H/Nu_0$, not $Nv_1Nv_2$. Apply the positive norm~\eqref{phy:input}
inside the exact common-divisor identity~\eqref{sw:mobius}.
For each squarefree good $d$, put $L=D/Nd$ and keep the original
coefficient $A_U(dn)$, the column mask $(n,dS)=1$, and the raw row mask
$(d,u)=1$. The row mask may be enlarged only inside the positive bound.
The assignment bound and ideal count give coefficient squared norm
$\ll_\epsilon B_W^2D^\epsilon L$ for $L\ge1$.
For $1/2\le L<1$ the only possible good residual is the unit; its direct
row count is bounded by the same polynomial up to an absolute factor.
There are no active columns for $L<1/2$.
Thus, only now taking the triangle over $d$,
\[
 |S_1^{AA}|\ll B_W^2(HD)^\epsilon
     \sum_{Nd\le2D}\frac1{Nd}\Phi_6(H/Nu_0,D/Nd).
\]
The four ideal sums have exponents $1,2,4/3,11/6$.
The first is logarithmic, and the others converge, which proves the
fixed-$u_0$ assertion of Proposition~\ref{phy:signed}.
If all $S$-supported $u_0$ are included, the terms with positive
$Nu_0$ exponent sum geometrically. The $D$ term has exponent zero
and costs only the finite count
$O_S((\log(2H))^{|S|})$, absorbed into $(HD)^\epsilon$.
It is not summed as a convergent exponent-zero Euler product.

The accepted identity of Section~\ref{cmr:core} gives
\[
 S_1^{AA}=\frac{K_0-K_{\rm all}}D+\mathcal E_{\rm corr},
 \qquad |\mathcal E_{\rm corr}|\le236800B_W^2R_w
\]
in its stated high window and at $H\le D^{3/4}$.
This retains its ordered $(p,q)$ correction, including $p=q$ and the
asymmetric restored prime mask; it is distinct from the unordered
unit correction $C_2$. Taking this already bounded correction gives
\eqref{phy:combined}. Neither $K_0$ nor $K_{\rm all}$ is separately
bounded by a substituted marginal mean. The resulting $13/12$ bound
is an improved positive comparison, not the target signed estimate.

\subsection{The updated short certificate and its unresolved input}\label{phy:short-proof}
For a nonempty block $H_b\ge1$. At $1\le L\le D$, the two already established block bounds and the new
physical-row bound are all valid positive estimates of the SAME actual
unweighted operator. Their minimum is Proposition~\ref{phy:short}.
The factor $8$ in $H_b$ is the literal dyadic bound
$Nv_1(Nv_2)^2<8A_1A_2^2$, with the original capacity retained.
No radical conductor is substituted for that physical norm.
At $1/2\le L<1$ retain the actual unit-column norm.
For the exact mean-retaining Abel identity of Proposition~\ref{cmr:joint},
write $X_0=D/U$ and $M_1(W)=\max(\|W\|_\infty,\|xW'\|_\infty)$.
Its sufficient positive discrepancy certificate is
\[
 \frac{|\mathcal E_{\rm PIT}|}{D}
 \ll (HD)^\epsilon M_1(W)^2\delta_S(U)
 \left\{\lambda_{{\rm now},*}(X_0)
       +\int_{1/2}^{X_0}\lambda_{{\rm now},*}(L)\frac{dL}{L}\right\}.
\]
Here $U\ge7$ is retained for the Li/Stieltjes definition, and the lower
endpoint $E_S(U)Q(U)$ is still part of the exact identity.
The minimum must stay inside the integral for general neighboring
capacity-feasible blocks. It is not a variance-only substitution.
At the critical line, the four new short monomial powers are
$3/4,1/2,19/24,2/3$. The pure minimum is $19/24$; in the mixed
$A_1\asymp D^{5/12},A_2\asymp D^{1/6}$ neighborhood the old $7/9$
entry remains smaller. Equating those majorant powers with the $3/4$
target gives the sufficient scalar exponents in~\eqref{phy:diagnostic}.
These are not actual-error lower bounds.
The conditional strip implication in Theorem~\ref{cmr:conditional-strip}
therefore has boundaries $11/12$ and $17/18$ for these proposed inputs;
no such input or strip has been established here. In particular,
pointwise endpoint control and the coherent signed $K_0-J$ are not
supplied by an independent mean-square assumption alone.

The new primary input improves a positive norm certificate. Equal-norm
ideal orientations, coherent means, original common profiles, full-union
raw zeros and all fixed-conductor qualifications remain. The actual
signed target, independent determinant bridge and analytic/category
Lean work remain open.
